\chapter{Lua-Skripting}
\label{ch:basics/lua}

In diesem Kapitel lernst du die Skripterstellung in Lua kennen, die
Werkzeuge, die dir dabei helfen, und einige Techniken, die du nützlich
finden wirst.

\section{Programmierung}\label{programmierung}

Beim Programmieren geht es darum, ein Problem -- etwa das Sortieren
einer Liste -- in Schritte zu zerlegen, die ein Computer verstehen kann.
Den logischen Prozess des Programmierens zu vermitteln, würde den Rahmen
dieses Buches sprengen; die folgenden Angebote sind dafür aber sehr
hilfreich:

\begin{itemize}
\tightlist
\item
  Codecademy\footnote{\url{http://www.codecademy.com/}} (englisch) ist
  eine der besten Quellen, um Programmieren zu lernen. Sie bietet
  interaktive Lerneinheiten.
\item
  Scratch\footnote{\url{https://scratch.mit.edu}} eignet sich gut, um
  bei den absoluten Grundlagen anzufangen und die
  Problemlösungstechniken zu erlernen, die man zum Programmieren
  braucht. Besonders gut geeignet für Kinder und Jugendliche.
\item
  The Morpheus\footnote{\url{https://www.youtube.com/@TheMorpheusTutorials}}
  ist ein guter YouTube-Kanal, um programmieren zu lernen.
\end{itemize}

\subsection{Programmieren in Lua}\label{programmieren-in-lua}

Auch Lua selbst zu lehren, würde über dieses Buch hinausgehen. Das Buch
Programming in Lua (PiL)\footnote{\url{https://www.lua.org/pil/contents.html}}
(englisch) ist jedoch eine ausgezeichnete Einführung in die
Programmierung mit Lua.

\section{Code-Editoren}\label{code-editoren}

Zum Schreiben von Lua-Skripten genügt ein Code-Editor mit
Syntaxhervorhebung. Dabei werden Wörter und Zeichen je nach Bedeutung
unterschiedlich eingefärbt. So fallen Fehler und Ungereimtheiten leicht
auf.

Zum Beispiel:

\begin{Shaded}
\begin{Highlighting}[]
\KeywordTok{function} \VariableTok{ctf}\OperatorTok{.}\NormalTok{post}\OperatorTok{(}\VariableTok{team}\OperatorTok{,} \VariableTok{msg}\OperatorTok{)}
    \ControlFlowTok{if} \KeywordTok{not} \VariableTok{ctf}\OperatorTok{.}\NormalTok{team}\OperatorTok{(}\VariableTok{team}\OperatorTok{)} \ControlFlowTok{then}
        \ControlFlowTok{return} \KeywordTok{false}
    \ControlFlowTok{end}
    \ControlFlowTok{if} \KeywordTok{not} \VariableTok{ctf}\OperatorTok{.}\NormalTok{team}\OperatorTok{(}\VariableTok{team}\OperatorTok{).}\VariableTok{log} \ControlFlowTok{then}
        \VariableTok{ctf}\OperatorTok{.}\NormalTok{team}\OperatorTok{(}\VariableTok{team}\OperatorTok{).}\VariableTok{log} \OperatorTok{=} \OperatorTok{\{\}}
    \ControlFlowTok{end}

    \FunctionTok{table.insert}\OperatorTok{(}\VariableTok{ctf}\OperatorTok{.}\NormalTok{team}\OperatorTok{(}\VariableTok{team}\OperatorTok{).}\VariableTok{log}\OperatorTok{,} \DecValTok{1}\OperatorTok{,} \VariableTok{msg}\OperatorTok{)}
    \VariableTok{ctf}\OperatorTok{.}\NormalTok{save}\OperatorTok{()}

    \ControlFlowTok{return} \KeywordTok{true}
\KeywordTok{end}
\end{Highlighting}
\end{Shaded}

In diesem Beispiel sind die Schlüsselwörter hervorgehoben, darunter
\texttt{if}, \texttt{then}, \texttt{end} und \texttt{return}.
Funktionen, die Lua standardmäßig mitbringt, etwa
\texttt{table.\allowbreak{}insert}, werden ebenfalls hervorgehoben.

Für den Einstieg empfehlen sich diese Editoren:

\begin{itemize}
\tightlist
\item
  VSCode\footnote{\url{https://code.visualstudio.com/}} bzw. die
  quelloffene Variante VSCodium\footnote{\url{https://vscodium.com/}}:
  verbreitet und für alle Betriebssysteme erhältlich. Installiere dazu
  die Erweiterung Lua\footnote{\url{https://marketplace.visualstudio.com/items?itemName=sumneko.lua}}
  (Lua Language Server). Sie vervollständigt Code automatisch, zeigt
  Fehler schon beim Tippen an und springt auf Wunsch zur Definition
  einer Funktion. Ergänzend gibt es Erweiterungen speziell für
  Luanti\footnote{\url{https://marketplace.visualstudio.com/items?itemName=GreenXenith.minetest-tools}}.
\item
  Notepad++\footnote{\url{http://notepad-plus-plus.org/}} (englisch):
  schlicht, nur für Windows.
\end{itemize}

Es gibt auch andere geeignete Editoren.

\section{Lokaler und globaler
Gültigkeitsbereich}\label{lokaler-und-globaler-guxfcltigkeitsbereich}

Ob eine Variable lokal oder global ist, bestimmt, von wo aus sie gelesen
und beschrieben werden kann. Auf globale Variablen kann von überall in
der Skriptdatei und auch von jedem anderen Mod aus zugegriffen werden:

\begin{Shaded}
\begin{Highlighting}[]
\KeywordTok{function}\NormalTok{ one}\OperatorTok{()}
    \VariableTok{foo} \OperatorTok{=} \StringTok{"bar"}
\KeywordTok{end}

\KeywordTok{function}\NormalTok{ two}\OperatorTok{()}
    \FunctionTok{print}\OperatorTok{(}\FunctionTok{dump}\OperatorTok{(}\VariableTok{foo}\OperatorTok{))} \CommentTok{{-}{-} Ausgabe: "bar"}
\KeywordTok{end}

\NormalTok{one}\OperatorTok{()}
\NormalTok{two}\OperatorTok{()}
\end{Highlighting}
\end{Shaded}

Auf eine lokale Variable kann dagegen nur dort zugegriffen werden, wo
sie definiert ist. In Lua sind Variablen standardmäßig global; du musst
daher ausdrücklich das Schlüsselwort \texttt{local} verwenden:

\begin{Shaded}
\begin{Highlighting}[]
\CommentTok{{-}{-} Innerhalb dieser Skriptdatei zugänglich}
\KeywordTok{local} \VariableTok{one} \OperatorTok{=} \DecValTok{1}

\KeywordTok{function}\NormalTok{ myfunc}\OperatorTok{()}
    \CommentTok{{-}{-} Innerhalb dieser Funktion zugänglich}
    \KeywordTok{local} \VariableTok{two} \OperatorTok{=} \VariableTok{one} \OperatorTok{+} \VariableTok{one}

    \ControlFlowTok{if} \VariableTok{two} \OperatorTok{==} \VariableTok{one} \ControlFlowTok{then}
        \CommentTok{{-}{-} Innerhalb dieser if{-}Anweisung zugänglich}
        \KeywordTok{local} \VariableTok{three} \OperatorTok{=} \VariableTok{one} \OperatorTok{+} \VariableTok{two}
    \ControlFlowTok{end}
\KeywordTok{end}
\end{Highlighting}
\end{Shaded}

\subsection{Möglichst lokale Variablen
verwenden}\label{muxf6glichst-lokale-variablen-verwenden}

Wann immer möglich, sollten lokale Variablen verwendet werden. Ein Mod
sollte höchstens eine globale Variable anlegen, und zwar mit demselben
Namen wie der Mod. Weitere globale Variablen sind unsauberer
Programmierstil, und Luanti warnt davor:

\begin{Shaded}
\begin{Highlighting}[]
\NormalTok{Assignment to undeclared global \textquotesingle{}foo\textquotesingle{} inside function}
\NormalTok{at init.lua:2}
\end{Highlighting}
\end{Shaded}

Sinngemäß: Zuweisung an die nicht deklarierte globale Variable
\texttt{foo} innerhalb einer Funktion in
\texttt{init.\allowbreak{}lua}\index{init.lua@\texttt{init.\allowbreak{}lua}},
Zeile 2.

Um das zu beheben, verwende \texttt{local}:

\begin{Shaded}
\begin{Highlighting}[]
\KeywordTok{function}\NormalTok{ one}\OperatorTok{()}
    \KeywordTok{local} \VariableTok{foo} \OperatorTok{=} \StringTok{"bar"}
\KeywordTok{end}

\KeywordTok{function}\NormalTok{ two}\OperatorTok{()}
    \FunctionTok{print}\OperatorTok{(}\FunctionTok{dump}\OperatorTok{(}\VariableTok{foo}\OperatorTok{))} \CommentTok{{-}{-} Ausgabe: nil}
\KeywordTok{end}

\NormalTok{one}\OperatorTok{()}
\NormalTok{two}\OperatorTok{()}
\end{Highlighting}
\end{Shaded}

Denk daran, dass \texttt{nil} \textbf{nicht initialisiert} bedeutet: Der
Variablen wurde noch kein Wert zugewiesen, sie existiert nicht, oder ihr
Wert wurde entfernt (also auf \texttt{nil} gesetzt).

Funktionen sind Variablen eines besonderen Typs, sollten aber ebenfalls
lokal sein, weil andere Mods Funktionen mit demselben Namen haben
könnten.

\begin{Shaded}
\begin{Highlighting}[]
\KeywordTok{local} \KeywordTok{function}\NormalTok{ foo}\OperatorTok{(}\VariableTok{bar}\OperatorTok{)}
    \ControlFlowTok{return} \VariableTok{bar} \OperatorTok{*} \DecValTok{2}
\KeywordTok{end}
\end{Highlighting}
\end{Shaded}

Damit andere Mods deine Funktionen aufrufen können, lege eine Tabelle
mit demselben Namen wie der Mod an und füge deine Funktionen dort ein.
Diese Tabelle wird oft als API-Tabelle oder
Namensraum\index{namensraum@Namensraum} (\emph{namespace}) bezeichnet.

\begin{Shaded}
\begin{Highlighting}[]
\VariableTok{mymod} \OperatorTok{=} \OperatorTok{\{\}}

\KeywordTok{function} \VariableTok{mymod}\OperatorTok{.}\NormalTok{foo}\OperatorTok{(}\VariableTok{bar}\OperatorTok{)}
    \ControlFlowTok{return} \StringTok{"foo"} \OperatorTok{..} \VariableTok{bar}
\KeywordTok{end}

\CommentTok{{-}{-} In einem anderen Mod oder Skript:}
\VariableTok{mymod}\OperatorTok{.}\NormalTok{foo}\OperatorTok{(}\StringTok{"foobar"}\OperatorTok{)}
\end{Highlighting}
\end{Shaded}

Die Schreibweise \texttt{function mymod.\allowbreak{}foo()} ist
gleichbedeutend mit \texttt{mymod.\allowbreak{}foo = function()}, nur
etwas lesbarer.

\section{Andere Lua-Skripte
einbinden}\label{andere-lua-skripte-einbinden}

Andere Lua-Skripte bindet man in einen Mod am besten mit
\emph{dofile\index{dofile@\texttt{dofile}}} ein.

\begin{Shaded}
\begin{Highlighting}[]
\FunctionTok{dofile}\OperatorTok{(}\VariableTok{core}\OperatorTok{.}\NormalTok{get\_modpath}\OperatorTok{(}\StringTok{"modname"}\OperatorTok{)} \OperatorTok{..} \StringTok{"/script.lua"}\OperatorTok{)}
\end{Highlighting}
\end{Shaded}

Ein Skript kann einen Wert zurückgeben. Damit lassen sich private lokale
Variablen gezielt weitergeben:

\begin{Shaded}
\begin{Highlighting}[]
\CommentTok{{-}{-} script.lua}
\KeywordTok{local} \VariableTok{module} \OperatorTok{=} \OperatorTok{\{\}}
\VariableTok{module}\OperatorTok{.}\VariableTok{message} \OperatorTok{=} \StringTok{"Hallo Welt!"}
\ControlFlowTok{return} \VariableTok{module}
\end{Highlighting}
\end{Shaded}

\begin{Shaded}
\begin{Highlighting}[]
\CommentTok{{-}{-} init.lua}
\KeywordTok{local} \VariableTok{ret} \OperatorTok{=}
    \FunctionTok{dofile}\OperatorTok{(}\VariableTok{core}\OperatorTok{.}\NormalTok{get\_modpath}\OperatorTok{(}\StringTok{"modname"}\OperatorTok{)} \OperatorTok{..} \StringTok{"/script.lua"}\OperatorTok{)}
\FunctionTok{print}\OperatorTok{(}\VariableTok{ret}\OperatorTok{.}\VariableTok{message}\OperatorTok{)} \CommentTok{{-}{-} Hallo Welt!}
\end{Highlighting}
\end{Shaded}

Wie man den Code eines Mods am besten aufteilt, wird in einem späteren
Kapitel~(Kapitel~\ref{ch:quality/clean-arch}) erklärt.
